%% =====================================================================
%%  Talleres GARyC -- Análisis Económico y Computacional
%%  Recursos y requerimientos
%%  Compilar con pdfLaTeX (Overleaf: Menú > Compilador > pdfLaTeX).
%% =====================================================================
\documentclass[11pt]{article}
\usepackage{garyc}

\documento{Recursos y requerimientos}
\title{Talleres GARyC: Análisis Económico y Computacional. Recursos y requerimientos}
\author{Talleres GARyC}
\date{}
\seccionesnumeradas

\begin{document}

\portada{Recursos y requerimientos}

\begin{presentacion}
Se presenta la relación de recursos materiales, tecnológicos y digitales
necesarios para impartir los tres talleres del syllabus. Su carácter
principalmente práctico requiere realizar ejercicios en vivo con R, Python o
Stata, además de integrar herramientas de IA generativa en el flujo de
trabajo.

Se propone trabajar en dos bloques de dos horas por semana (cuatro horas
semanales). Esta distribución permitiría cubrir las 53 horas previstas en
aproximadamente 14 semanas, agrupando hasta dos sesiones del syllabus por
bloque.
\end{presentacion}

%% ---------------------------------------------------------------------
\section{Infraestructura, espacios y equipo físico}
\begin{recursos}
\recurso{Sala de juntas} Reserva de la sala de juntas del piso~5 del
  edificio Guardiola en dos bloques de dos horas semanales, durante todo el
  programa.
\recurso{Laptops para facilitadores} Dos laptops institucionales del Banco,
  con cargadores y permisos de uso e instalación del software requerido.
\recurso{Proyección y conexión} Un cable HDMI o los adaptadores que requieran
  los equipos para conectarse a la pantalla de la sala de juntas y acceso a
  red y a tomas de corriente.
\recurso{Material auxiliar} Extensiones o multicontactos.
\end{recursos}

%% ---------------------------------------------------------------------
\section{Software y recursos digitales}
\begin{recursos}
\recurso{Stata} Versión institucional más reciente disponible y licencias
  vigentes para los equipos que lo utilicen.
\recurso{R, RStudio y Rtools} Versiones institucionalmente disponibles y
  compatibles entre sí.
\recurso{Python y Anaconda} Python, Anaconda y entornos de trabajo para
  ejecutar scripts y cuadernos de Jupyter.
\recurso{Bibliotecas y complementos} Paquetes para econometría causal, datos
  panel, aprendizaje automático y análisis espacial; dependencias instaladas
  y probadas de antemano.
\recurso{Accesos y permisos} Instalación autorizada y, en su caso, acceso a
  CRAN, PyPI/conda y SSC.
\recurso{Repositorio compartido} Carpeta institucional con acceso para
  participantes, organizada por taller, módulo y sesión, para distribuir
  materiales y entregables.
\recurso{Herramienta de IA generativa} Acceso a una solución autorizada
  institucionalmente.
\end{recursos}

\end{document}
